\documentclass[]{article}
\usepackage{lmodern}
\usepackage{amssymb,amsmath}
\usepackage{ifxetex,ifluatex}
\usepackage{fixltx2e} % provides \textsubscript
\ifnum 0\ifxetex 1\fi\ifluatex 1\fi=0 % if pdftex
  \usepackage[T1]{fontenc}
  \usepackage[utf8]{inputenc}
\else % if luatex or xelatex
  \ifxetex
    \usepackage{mathspec}
  \else
    \usepackage{fontspec}
  \fi
  \defaultfontfeatures{Ligatures=TeX,Scale=MatchLowercase}
\fi
% use upquote if available, for straight quotes in verbatim environments
\IfFileExists{upquote.sty}{\usepackage{upquote}}{}
% use microtype if available
\IfFileExists{microtype.sty}{%
\usepackage{microtype}
\UseMicrotypeSet[protrusion]{basicmath} % disable protrusion for tt fonts
}{}
\usepackage[margin=1in]{geometry}
\usepackage{hyperref}
\hypersetup{unicode=true,
            pdftitle={ACTIVIDAD 1. PIRHI},
            pdfborder={0 0 0},
            breaklinks=true}
\urlstyle{same}  % don't use monospace font for urls
\usepackage{graphicx,grffile}
\makeatletter
\def\maxwidth{\ifdim\Gin@nat@width>\linewidth\linewidth\else\Gin@nat@width\fi}
\def\maxheight{\ifdim\Gin@nat@height>\textheight\textheight\else\Gin@nat@height\fi}
\makeatother
% Scale images if necessary, so that they will not overflow the page
% margins by default, and it is still possible to overwrite the defaults
% using explicit options in \includegraphics[width, height, ...]{}
\setkeys{Gin}{width=\maxwidth,height=\maxheight,keepaspectratio}
\IfFileExists{parskip.sty}{%
\usepackage{parskip}
}{% else
\setlength{\parindent}{0pt}
\setlength{\parskip}{6pt plus 2pt minus 1pt}
}
\setlength{\emergencystretch}{3em}  % prevent overfull lines
\providecommand{\tightlist}{%
  \setlength{\itemsep}{0pt}\setlength{\parskip}{0pt}}
\setcounter{secnumdepth}{5}

%%% Use protect on footnotes to avoid problems with footnotes in titles
\let\rmarkdownfootnote\footnote%
\def\footnote{\protect\rmarkdownfootnote}

%%% Change title format to be more compact
\usepackage{titling}

% Create subtitle command for use in maketitle
\newcommand{\subtitle}[1]{
  \posttitle{
    \begin{center}\large#1\end{center}
    }
}

\setlength{\droptitle}{-2em}
  \title{Diagnóstico técnico-social del protocolo de priorización de inversiones en mejoramiento de la eficiencia de conducción hídrica.}
  \pretitle{\vspace{\droptitle}\centering\huge}
  \posttitle{\par}
  \author{}
  \preauthor{}\postauthor{}
  \date{}
  \predate{}\postdate{}

\usepackage{geometry}
\geometry{letterpaper}
\usepackage{graphicx}
\usepackage{amssymb}
\usepackage{hyperref}
\usepackage{titlesec}
\usepackage{appendix}  
\usepackage{booktabs}
\usepackage{longtable}
\usepackage{multirow}
\usepackage{subcaption}
\usepackage{rotating}

\usepackage{fontspec,xltxtra,xunicode}
\defaultfontfeatures{Mapping=tex-text}
\setmainfont[		
 BoldFont={Century Gothic Bold}, 
 ItalicFont={Century Gothic Italic},
 BoldItalicFont={Century Gothic Bold Italic}
 ]{Century Gothic}

\renewcommand{\figurename}{Figura}
\renewcommand{\tablename}{Cuadro} 
\renewcommand{\contentsname}{Tabla de contenidos\\} 
\renewcommand{\listfigurename}{Índice de figuras\\}
\renewcommand{\listtablename}{Índice de cuadros\\} 
\renewcommand{\appendixname}{Anexos}
\renewcommand{\appendixtocname}{Anexos}
\renewcommand{\appendixpagename}{Anexos}


\titleformat*{\section}{\normalsize\bfseries}
\titleformat*{\subsection}{\normalsize\bfseries}
\titleformat*{\subsubsection}{\normalsize\bfseries}

\usepackage{fancyvrb}
\newcommand{\VerbBar}{|}
\newcommand{\VERB}{\Verb[commandchars=\\\{\}]}
\DefineVerbatimEnvironment{Highlighting}{Verbatim}{commandchars=\\\{\}}
% Add ',fontsize=\small' for more characters per line
\usepackage{color}
\usepackage{framed}
\usepackage{tabulary}
\definecolor{shadecolor}{RGB}{248,248,248}
\newenvironment{Shaded}{\begin{snugshade}}{\end{snugshade}}
\newcommand{\KeywordTok}[1]{\textcolor[rgb]{0.13,0.29,0.53}{\textbf{#1}}}
\newcommand{\DataTypeTok}[1]{\textcolor[rgb]{0.13,0.29,0.53}{#1}}
\newcommand{\DecValTok}[1]{\textcolor[rgb]{0.00,0.00,0.81}{#1}}
\newcommand{\BaseNTok}[1]{\textcolor[rgb]{0.00,0.00,0.81}{#1}}
\newcommand{\FloatTok}[1]{\textcolor[rgb]{0.00,0.00,0.81}{#1}}
\newcommand{\ConstantTok}[1]{\textcolor[rgb]{0.00,0.00,0.00}{#1}}
\newcommand{\CharTok}[1]{\textcolor[rgb]{0.31,0.60,0.02}{#1}}
\newcommand{\SpecialCharTok}[1]{\textcolor[rgb]{0.00,0.00,0.00}{#1}}
\newcommand{\StringTok}[1]{\textcolor[rgb]{0.31,0.60,0.02}{#1}}
\newcommand{\VerbatimStringTok}[1]{\textcolor[rgb]{0.31,0.60,0.02}{#1}}
\newcommand{\SpecialStringTok}[1]{\textcolor[rgb]{0.31,0.60,0.02}{#1}}
\newcommand{\ImportTok}[1]{#1}
\newcommand{\CommentTok}[1]{\textcolor[rgb]{0.56,0.35,0.01}{\textit{#1}}}
\newcommand{\DocumentationTok}[1]{\textcolor[rgb]{0.56,0.35,0.01}{\textbf{\textit{#1}}}}
\newcommand{\AnnotationTok}[1]{\textcolor[rgb]{0.56,0.35,0.01}{\textbf{\textit{#1}}}}
\newcommand{\CommentVarTok}[1]{\textcolor[rgb]{0.56,0.35,0.01}{\textbf{\textit{#1}}}}
\newcommand{\OtherTok}[1]{\textcolor[rgb]{0.56,0.35,0.01}{#1}}
\newcommand{\FunctionTok}[1]{\textcolor[rgb]{0.00,0.00,0.00}{#1}}
\newcommand{\VariableTok}[1]{\textcolor[rgb]{0.00,0.00,0.00}{#1}}
\newcommand{\ControlFlowTok}[1]{\textcolor[rgb]{0.13,0.29,0.53}{\textbf{#1}}}
\newcommand{\OperatorTok}[1]{\textcolor[rgb]{0.81,0.36,0.00}{\textbf{#1}}}
\newcommand{\BuiltInTok}[1]{#1}
\newcommand{\ExtensionTok}[1]{#1}
\newcommand{\PreprocessorTok}[1]{\textcolor[rgb]{0.56,0.35,0.01}{\textit{#1}}}
\newcommand{\AttributeTok}[1]{\textcolor[rgb]{0.77,0.63,0.00}{#1}}
\newcommand{\RegionMarkerTok}[1]{#1}
\newcommand{\InformationTok}[1]{\textcolor[rgb]{0.56,0.35,0.01}{\textbf{\textit{#1}}}}
\newcommand{\WarningTok}[1]{\textcolor[rgb]{0.56,0.35,0.01}{\textbf{\textit{#1}}}}
\newcommand{\AlertTok}[1]{\textcolor[rgb]{0.94,0.16,0.16}{#1}}
\newcommand{\ErrorTok}[1]{\textcolor[rgb]{0.64,0.00,0.00}{\textbf{#1}}}
\newcommand{\NormalTok}[1]{#1}
\usepackage{graphicx,grffile}
\makeatletter
\def\maxwidth{\ifdim\Gin@nat@width>\linewidth\linewidth\else\Gin@nat@width\fi}
\def\maxheight{\ifdim\Gin@nat@height>\textheight\textheight\else\Gin@nat@height\fi}
\makeatother
\usepackage{booktabs}
\usepackage{longtable}
\usepackage{array}
\usepackage{multirow}
\usepackage[table]{xcolor}
\usepackage{wrapfig}
\usepackage{float}
\usepackage{colortbl}
\usepackage{pdflscape}
\usepackage{tabu}
\usepackage{threeparttable}
\usepackage[normalem]{ulem}
\usepackage[none]{hyphenat}
\usepackage{multicol}

% Regla con color
\makeatletter
\let\old@rule\@rule
\def\@rule[#1]#2#3{\textcolor{rulecolor}{\old@rule[#1]{#2}{#3}}}
\makeatother
% Cambiar color de regla
\definecolor{rulecolor}{RGB}{183,183,183}

% Color para las letras del título
\definecolor{casinegro}{RGB}{71,71,71}

\begin{document}
\sloppy 

\newgeometry{left=3cm,bottom=0cm,top=1cm,right=3cm}

\begin{titlepage}

\begin{figure} [H]
  \includegraphics[width=.2\textwidth]{Logo/bg_gob_logo.png}
\end{figure}

\fontsize{15}{0} \selectfont{\textbf{\textcolor{casinegro}{``PLANIFICACIÓN DE LA INVERSIÓN EN REVESTIMIENTOS HÍDRICOS"}}} \vspace*{.1in}\\ % el inicio de comillas simples es ` y de doble es de ``.
\fontsize{15}{0} \selectfont{\textcolor{casinegro}{FIC-R 2016}}\\
\rule{10cm}{.05cm}
\vspace*{.1in}
\hspace*{.2in}

\fontsize{14}{0} \selectfont{\textcolor{casinegro}{Universidad de La Serena}} \vspace*{.1in}\\
\fontsize{14}{0} \selectfont{\textbf{\textcolor{casinegro}{Región de Coquimbo}}}

\vspace*{.1in}

\begin{figure}[H]
  \centering
\begin{subfigure}[c]{.2\textwidth}
  \includegraphics[width=\textwidth]{Logo/GOB_RGB.png}
\end{subfigure}
\hfill
\begin{subfigure}[c]{.23\textwidth}
  \includegraphics[width=\textwidth]{Logo/uls.png} 
\end{subfigure}
\end{figure}
\vspace*{4in}

\begin{figure}[H]
  \centering
\begin{subfigure}{.74\textwidth}
   \fontsize{14}{0} \selectfont{\textcolor{casinegro}{Marzo 2018}} \\
   {\definecolor{rulecolor}{named}{casinegro}\rule[-3mm]{10cm}{.05cm}}
\end{subfigure}
\hfill
\begin{subfigure}{.25\textwidth}
\hfill
  \includegraphics[width=\textwidth]{Logo/lgore.pdf} 
\end{subfigure}
\end{figure}


\end{titlepage}

\restoregeometry


\begin{table}[!htb]
\centering \vspace*{2.5in}
\begin{tabular}{|l|l|l|l|}
    \hline
    \multicolumn{3}{|c|}{\multirow{2}{8cm}{\centering \textbf{INFORME DE AVANCE TÉCNICO} \\
    \textbf{FIC REGIONAL} }} & \textbf{Fecha de emisión:}\\
    \cline{4-4}
    \multicolumn{3}{|c|}{\multirow{2}{8cm}} &  Abril 2018\\
    \hline
    \textbf{Fecha de inicio:}  &\textbf{Fecha de término:} &\textbf{N° de informe:} &\textbf{Duración:}\\
     20/Diciembre/2017 & 20/Octubre/2019 & 1 & 22 meses\\
    \hline
    \textbf{Nombre del proyecto:} & \multicolumn{3}{|l|}{\textbf{Planificación de la Inversión en Recursos Hídricos}}\\
    \hline
    \multirow{4}{4cm}{\textbf{Financiamiento:}} & \textbf{Monto total del proyecto:} & \multicolumn{2}{c|}{\centering M\$ 222.188}\\
    \cline{2-4}
    & \multirow{2}{4cm}{\textbf{Aporte FIC-R:}} & \multicolumn{2}{c|}{\textbf{Año 1:} M\$ 90.040}\\
    \cline{3-4}
    & & \multicolumn{2}{c|}{\textbf{Año 2:} M\$ 109.928}\\
    \cline{2-4}
    & \textbf{Aporte institución:} & \multicolumn{2}{c|}{\centering M\$ 22.220}\\
    \hline
    \textbf{Institución ejecutora:} & \multicolumn{3}{|l|}{\centering Universidad de La Serena}\\
    \hline
    \textbf{Responsable del proyecto:} & \multicolumn{3}{|l|}{\centering Dr. Pablo Álvarez Latorre}\\
    \hline
    \textbf{Cargo:} & \multicolumn{3}{|l|}{\centering Director del proyecto}\\
    \hline
    \textbf{E-mail:} & \multicolumn{3}{|l|}{\centering pabloa@userena.cl}\\
    \hline
    \textbf{Fono:} & \multicolumn{3}{|l|}{\centering 51-2554913}\\
    \hline

\end{tabular}

\end{table}

\clearpage

\section{Resumen del Proyecto}

La iniciativa denominada ``Planificación de la Inversión en Revestimientos Hídricos", se adjudicó un financiamiento FIC de \$ 199.968.000, lo que sumado a los aportes de la contraparte totalizan \$ 220.188.000. Ejecutándose desde Diciembre de 2017 a Octubre de 2019, por un período de 22 meses.\\
\\
El proyecto tiene por finalidad planificar el mejoramiento de la eficiencia de conducción hídrica, a través de la transferencia de conocimientos y herramientas a las Organizaciones de Usuarios de Aguas para apoyar la gestión de recursos hídricos. El impacto directo o propósito es incrementar la oferta neta, por medio de la priorización de las inversiones público-privadas basados en parámetros técnico-social.\\
\\
El área en la cual se desarrolla esta iniciativa corresponde a la cuenca del río Choapa, abarcando toda el área de influencia de la Junta de Vigilancia del río Chalinga y sus Afluentes, Junta de Vigilancia del río Illapel y sus Afluentes, y la Junta de Vigilancia del río Choapa y sus Afluentes. Por lo mismo, las instituciones asociadas al proyecto fueron la Junta de Vigilancia río Elqui y sus afluentes, Junta de Vigilancia río Illapel y Junta de Vigilancia Estero Derecho. Por lo mismo, las instituciones asociadas al pryectos son la Junta de Vigilancia del río Chalinga y sus Afluentes, la Junta de Vigilancia del río Illapel y sus afluentes, y la Junta de Vigilancia del río Choapa y sus Afluentes, además como beneficiario indirecto la Fundación Minera Los Pelambres.\\
\\
Los resultados esperados de este proyecto son:\\
\begin{itemize}
\item Protocolo técnico-social de priorización de la inversión en mejoramiento de la eficiencia de conducción hídrica.
\item Plan de priorización de la inversión en mejoramiento de la eficiencia de conducción hídrica para la red de distribución de la Junta de Vigilancia Río Chalinga y sus Afluentes.
\item Plan de priorización de la inversión en mejoramiento de la eficiencia de conducción hídrica para la red de distribución de la Junta de Vigilancia Río Illapel y sus Afluentes.
\item Plan de priorización de la inversión en mejoramiento de la eficiencia de conducción hídrica para la red de distribución en la Junta de Vigilancia Río Choapa y sus Afluentes.
\item Programa de capacitación en gestión hídrica.
\item Sistema virtual de priorización de la inversión en mejoramiento de la eficiencia de conducción hídrica para las cuencas del río Choapa.
\item Protocolo técnico-social de priorización de la inversión en mejoramiento de la eficiencia de conducción hídrica.
\end{itemize}

El mérito innovador del proyecto se visualiza en establecer bases para la gestión integrada de los recursos hídricos de la cuenca, a través de un ordenamiento de la inversión en mejoramiento de eficiencia de conducción, por medio de la modelación dinámica del sistema hídrico.\\
\\
El equipo de trabajo de este proyecto está constituido por profesionales del Laboratorio de Prospección, Monitoreo y Modelación de Recursos Agrícolas y Ambientales (PROMMRA) perteneciente al Departamento de Agronomía de la Universidad de La Serena, liderado por el Dr. Pablo Álvarez Latorre.\\
\clearpage

\section{Actividades realizadas}

%\begin{table}[!htb]
%\centering
%\begin{tabular}{|c{2cm}|c{2cm}|c{2cm}|c{2cm}|c{2cm}|c{2cm}|}
 %   \hline
  %  \textbf{Componente} & \textbf{N° Actividad} & \textbf{Actividades programadas} & \textbf{Discrepancias} & \textbf{\% de avance físico}
   % \multirow {4}{2cm}{1. Protocolo técnico-social de Priorización de la Inversión en Mejorameinto de la Eficiencia de Conducción Hídrica.} & 		1 & 		Diagnóstico técnico-social del Protocolo de Priorización de Inversiones en Mejoramiento de la eficiencia de Conducción Hídrica & Ejecutada (E)      		& No & 100\%\\
%    \cline{2-6}
 %   & 1 & Diagnóstico técnico-sociañ del Protocolo de Priorización de Inversiones en Mejoramiento de la eficiencia de Conducción Hídrica & Ejecutada (E)      	& No & 100\%\\
 %   \cline{2-6}
  %  & 1 & Diagnóstico técnico-sociañ del Protocolo de Priorización de Inversiones en Mejoramiento de la eficiencia de Conducción Hídrica & Ejecutada (E)      	& No & 100\%\\
 %   \cline{2-6}
  %  & 1 & Diagnóstico técnico-sociañ del Protocolo de Priorización de Inversiones en Mejoramiento de la eficiencia de Conducción Hídrica & Ejecutada (E)      	& No & 100\%\\
 %   \cline{2-6}
  %  & 1 & Diagnóstico técnico-sociañ del Protocolo de Priorización de Inversiones en Mejoramiento de la eficiencia de Conducción Hídrica & Ejecutada (E)      	& No & 100\%\\
%    \cline{2-6}
 %   & 1 & Diagnóstico técnico-sociañ del Protocolo de Priorización de Inversiones en Mejoramiento de la eficiencia de Conducción Hídrica & Ejecutada (E)      	& No & 100\%\\
  %  \cline{2-6}
      
%    \multicolumn{3}{|c|}{\multirow{2}{8cm}{\centering \textbf{INFORME DE AVANCE TÉCNICO} \\
%    \textbf{FIC REGIONAL} }} & \textbf{Fecha de emisión:}\\
%    \cline{4-4}
%    \multicolumn{3}{|c|}{\multirow{2}{8cm}} &  Abril 2018\\
%    \hline
%    \textbf{Fecha de inicio:}  &\textbf{Fecha de término:} &\textbf{N° de informe:} &\textbf{Duración:}\\
%     20/Diciembre/2017 & 20/Octubre/2019 & 1 & 22 meses\\
%    \hline
%    \textbf{Nombre del proyecto:} & \multicolumn{3}{|l|}{\textbf{Planificación de la Inversión en Recursos Hídricos}}\\
%    \hline
%    \multirow{4}{4cm}{\textbf{Financiamiento:}} & \textbf{Monto total del proyecto:} & \multicolumn{2}{c|}{\centering M\$ 222.188}\\
%    \cline{2-4}
%    & \multirow{2}{4cm}{\textbf{Aporte FIC-R:}} & \multicolumn{2}{c|}{\textbf{Año 1:} M\$ 90.040}\\
%    \cline{3-4}
%    & & \multicolumn{2}{c|}{\textbf{Año 2:} M\$ 109.928}\\
%    \cline{2-4}
%    & \textbf{Aporte institución:} & \multicolumn{2}{c|}{\centering M\$ 22.220}\\
%    \hline
%    \textbf{Institución ejecutora:} & \multicolumn{3}{|l|}{\centering Universidad de La Serena}\\
%    \hline
%    \textbf{Responsable del proyecto:} & \multicolumn{3}{|l|}{\centering Dr. Pablo Álvarez Latorre}\\
%    \hline
%    \textbf{Cargo:} & \multicolumn{3}{|l|}{\centering Director del proyecto}\\
%    \hline
%    \textbf{E-mail:} & \multicolumn{3}{|l|}{\centering pabloa@userena.cl}\\
%    \hline
%    \textbf{Fono:} & \multicolumn{3}{|l|}{\centering 51-2554913}\\
%    \hline
%
%\end{tabular}
%
%\end{table}

\clearpage
\section{Carta Gantt}



\clearpage
\section{Metodología}

\subsection{Componente 1. Actividad 1.1. Diagnóstico técnico-social del Protocolo de Priorización de Inversiones en Mejoramiento de la eficiencia de Conducción Hídrica}

%incluir tabla

La priorización de la inversión público-privada en mejoramiento de la eficiencia de conducción hídrica, es primordial para resguardar el buen uso de los recursos disponibles, especialmente en una región donde el agua es un elemento vital para su desarrollo.

El proyecto FIC-R 2013 denominado \textit{``Diseño y puesta en marcha de un sistema de monitoreo de caudales y de un protocolo de determinación de pérdidas de agua, para la priorización de las inversiones público-privadas (SIMCA – ELQUI)"}, generó un protocolo de priorización de la inversión en la cuenca del río Elqui, para mejoramiento en la infraestructura de revestimiento de canales. El objetivo de este protocolo es unificar, estandarizar y normalizar los criterios, procedimientos y decisiones ligadas al mejoramiento de canales pertenecientes a una organización de regantes, basado en la variable técnica de pérdidas generales (l/s) por tramo y/o canal.

En la actualidad, el uso de los recursos hídricos se proyecta a nivel de cuenca, a través de un manejo integrado y sustentable de los diferentes usuarios de aguas. Es por esto, que se busca incluir variables de tipo social, para fortalecer la toma de decisiones en la distribución de recursos para inversión público-privada, y así realizar mejoramientos de la eficiencia de conducción hídrica en sectores donde se promueva el crecimiento y desarrollo, además de no interceder en la disponibilidad de los recursos de otras zonas.

Se realizó una revisión y diagnóstico del protocolo creado en SIMCA – ELQUI, con el fin de generar una propuesta de fortalecimiento que incluya las variables del ámbito social.

\subsubsection{Recopilación de antecedentes generales}

Se establecieron diferentes instancias para la revisión y diagnóstico del protocolo de priorización de inversiones en mejoramiento de la eficiencia de conducción hídrica. Se estableció una mesa técnica compuesta por los profesionales de las diferentes áreas del equipo, con el fin de identificar la información utilizada, metodología de recopilación y procesamiento de esta. Los resultados de esta revisión se presentan a continuación:

\paragraph{Recopilación antecedentes generales cuenca:}

Se realiza a través de una revisión bibliográfica de fuentes secundarias, como estudios de organismos referidos a los recursos hídricos de la cuenca a intervenir, como ejemplo se presentan las siguientes fuentes:

\begin{itemize}
\item Fichas de registro de bocatomas, Dirección General de Aguas
\item Canales de riego a nivel de matrices hasta 2do derivado, Comisión Nacional de Riego
\item Estudios de recursos hídricos de la zona a evaluar
\end{itemize}

\paragraph{Recopilación de antecedentes a nivel organizacional:}

Se realiza a través de reuniones y mesas de trabajo sostenidas con las organizaciones de usuarios de agua, donde se recoge información primaria la cual se detalla a continuación: 

\begin{itemize}
\item Diagrama unifilar
\item Ubicación geográfica bocatomas
\item Número de acciones totales
\item Número de acciones por comunidad de aguas
\item Dotaciones o desmarques
\item Capacidad de porteo canal
\item Número de usuarios comunidad de aguas
\item Contacto celadores
\end{itemize}

\subsubsection{Levantamiento de datos en terreno, cauce natural y artificial (canal).}

La caracterización del cauce natural, se realiza por medio de la georreferenciación y descripción de bocatomas, trazado de canales de aducción y compuertas de carga perteneciente a la comunidad de aguas. Con los puntos de trazado, se puede determinar la longitud del canal de aducción y su recorrido.

Posteriormente, en la caracterización del cauce artificial, se realiza la georreferenciación y descrpción de la infraestructura de conducción y distribución del recurso hídrico. Se capturan todas las obras de arte presentes en el recorrido del canal, desde la compuerta de carga hasta el final del recorrido, además se detalla el tipo de obra, material, finalidad, revestimeintos, además de puntos críticos como filtraciones o alteraciones visibles, quebradas, etc.

La caracterización de cauce natural y artificial se genera con el fin de proporcionar esta información a las OUA’s, como base de datos en formatos que permitan su visualización. Además con esta información se generan las variables que permitan realizar la determinación de pérdidas y posterior priorización de tramos y/o canales para el mejoramiento de la eficiencia de conducción hídrica.

\subsection{Determinación de pérdidas}

Con los insumos levantados en la caracterización, se procede a generar los tramos de aforo para la determinación de pérdidas. Estos tramos se dividen en matriz y distribución en primera instancia, para luego establecer los tramos por el cambios en el revestimiento presente, a través de aforos de los caudales de inicio y término de estos.







\subsection{Protocolo de priorización de inversiones en mejoramiento de la eficiencia de conducción hídrica}

Las variables que contiene el protocolo de priorización aplicado en SIMCA – ELQUI, las cuales se seleccionaron estadísticamente a través de la correlación de PEARSON, mayor a 0,80, en relación a la variable pérdidas generales (l/s) de los tramos.

\clearpage
\section{Resultados}

\end{document}